\documentclass[11pt]{article}
\usepackage[a4paper,margin=25mm]{geometry}
\usepackage{amsmath,amssymb}
\setlength{\parindent}{0pt}
\setlength{\parskip}{7pt}
\setlength{\emergencystretch}{2em}
\begin{document}
\begin{center}
{\Large Disproof of conjecture 00000006826}\\[5pt]
Hyperinvariant Schur triangularization\\
ziangni-sys
\end{center}
\textbf{Source and scope.}
The English source asserts an upper-triangular realization whose triangularization consists of hyperinvariant subspaces. The Chinese source likewise explicitly strengthens Schur triangularization to hyperinvariant subspaces. We use the standard meaning of a complete triangularizing flag: in complex dimension two it contains a one-dimensional intermediate subspace. A subspace is hyperinvariant when it is invariant under every bounded linear operator commuting with the given operator.

The counterexample refutes this strengthened requirement. Ordinary Schur triangularization remains valid. Indeed the example is already upper triangular in every basis, and ordinary invariant lines exist.

\textbf{The actual operator.}
Let $E=\mathbb C^2$ and $A=I_E$. Its matrix in the standard coordinate basis is
\[
 \begin{pmatrix}1&0\\0&1\end{pmatrix}.
\]
This is a bounded complex linear operator. Every bounded complex linear operator $S$ on $E$ commutes with $A$, since $SA=S=AS$.

\textbf{Classification of hyperinvariant subspaces.}
Suppose $M$ is a hyperinvariant linear subspace of $E$ and $M\ne\{0\}$. Choose $v\in M$ with $v\ne0$, and a coordinate $j\in\{0,1\}$ with $v_j\ne0$. For an arbitrary $w\in E$, define
\[
 S_w(z)=z_j(v_j)^{-1}w .
\]
Coordinate evaluation and multiplication by a fixed vector are continuous linear maps, so $S_w$ is an actual bounded complex linear operator. It commutes with $A$. Moreover,
\[
 S_w(v)=v_j(v_j)^{-1}w=w .
\]
Hyperinvariance gives $w\in M$. Since $w$ was arbitrary, $M=E$. Thus every hyperinvariant subspace of $I_E$ is either $\{0\}$ or $E$.

In particular none has complex dimension one. A complete triangularizing flag in dimension two must have the form
\[
 \{0\}\subset M_1\subset E,\qquad \dim_{\mathbb C}M_1=1.
\]
It cannot consist of hyperinvariant subspaces. This disproves the stated hyperinvariant strengthening.
\newpage
\textbf{Ordinary invariant subspaces are different.}
The coordinate line
\[
 L=\{(z_0,z_1)\in\mathbb C^2:z_1=0\}
\]
is invariant under $A=I_E$. It is nonzero because $(1,0)\in L$, and proper because $(1,1)\notin L$. It is not hyperinvariant: the classification just proved excludes all nonzero proper subspaces. Thus an ordinary invariant flag does not repair the hyperinvariant assertion.

The proof applies even to arbitrary algebraic subspaces, without a closedness assumption. Consequently it also applies to the usual convention that hyperinvariant subspaces are closed. In the present finite-dimensional space every linear subspace is closed anyway.

\textbf{Formal correspondence.}
The Lean package uses the actual complex vector space \texttt{Fin 2} $\to\mathbb C$, its continuous linear identity, and actual submodules. Hyperinvariance is defined by quantifying over continuous complex linear maps with the equality of composed operators. The rank-one map is constructed by composing coordinate evaluation with scalar multiplication by a fixed vector, through Mathlib's continuous linear map operations.

The formal proof shows that all maps commute with the identity, that the rank-one map sends the chosen nonzero vector to an arbitrary target, and that every hyperinvariant submodule is bottom or top. Actual complex finite dimension is two; hence no hyperinvariant submodule has dimension one. It also constructs the ordinary coordinate line as a genuine linear-map kernel, proves it nonzero and proper, and proves it is not hyperinvariant.

A separate actual identity matrix acts as the continuous linear identity on every vector. Every entry below its diagonal is zero. The final theorem combines this genuine upper-triangular matrix with the nonexistence of a complete hyperinvariant flag. No scalar proxy for invariant-subspace structure is assumed.

\textbf{Reproduction.}
The submission contains \texttt{lean/Main.lean}, a pinned public Lean project, this complete source, and the compiled PDF. With Lean 4.19.0 installed, run \texttt{lake build} inside the \texttt{lean} directory. Mathlib is pinned to commit
\[
\texttt{c44e0c8ee63ca166450922a373c7409c5d26b00b}.
\]
Printed axiom audits cover the commuting maps, rank-one action, classification, actual dimension, obstruction to an intermediate subspace, ordinary line, matrix action, triangularity, and final counterexample. The proof uses no admitted statements or extra axioms.

\textbf{Conclusion.}
The complex identity operator has an upper-triangular realization, but has no complete flag of hyperinvariant subspaces. The stronger Schur-type assertion therefore fails even in dimension two.
\end{document}
